\subsection{Coherent Frobenius: the collision obstruction and an actual normalization repair}
Here Frobenius means absolute Frobenius on characteristic-$p$ coherent
objects. It is not yet geometric Frobenius on an $\ell$-adic Weil
realization. The distinction between pullback on the singular center
and pullback on the parameter trait gives different, computable answers.

\begin{finding}[\fid{F36}: unit Frobenius depends on the coherent base]\label{f:coherentfrobenius}
Let $k=\F_q$ have characteristic $p>3$, let $R=k[[t]]$, and put
\[
 \Delta=t^m u(t),\quad u(t)\in R^*,\quad m\ge1,\qquad
 C=R[x]/(x(x-\Delta)),\quad L_1=C/(x),\quad L_2=C/(x-\Delta).
\]
For absolute Frobenius $F_C(c)=c^p$,
\[
 F_C^*L_i\simeq R\oplus R/(\Delta^{p-1})
 \quad\text{as }R\text{-modules}.
\]
Thus no unit $C$-linear Frobenius isomorphism $F_C^*L_i\simeq L_i$
exists. Every positive Tor degree of the derived pullback is also
nonzero. In the actual algebra and 64-dimensional module of
\fid{F19}/\fid{F21}, the center has $m=2$ and underived algebra
pullback adds torsion of $R$-length $80(p-1)$.

Nevertheless the actual module has a unit \emph{parameter-base}
semilinear Frobenius structure, compatible with the intrinsic
entrywise Frobenius on $A=R\times R\times M_2(C)$ after transport
through its regular adapted basis. On the normalized center the
actual module has a unit coherent Frobenius model preserving its
original $A$-action under restriction. Pushing this model back through
the collision gives the nonunit surjection just computed.
\end{finding}
\begin{proof}
\par\smallskip\noindent\textbf{The two branches.}\quad
Since $x^2=\Delta x$, characteristic $p$ gives
\[
 x^p=\Delta^{p-1}x,\qquad
 (x-\Delta)^p=\Delta^{p-1}(x-\Delta).
\]
Tensoring the cyclic presentations along $F_C$ gives
\[
 F_C^*L_1=C/(\Delta^{p-1}x),\qquad
 F_C^*L_2=C/(\Delta^{p-1}(x-\Delta)).
\]
The ring $C$ is free over $R$ with basis $1,x$. The first ideal is
exactly $R\Delta^{p-1}x$, because
$x(a+bx)=(a+b\Delta)x$. This proves the stated $R$-module
decomposition for $L_1$; use the basis $1,x-\Delta$ for $L_2$.
The torsion has length $m(p-1)$, while each $L_i$ is $R$-free.
Hence no possible $C$-linear isomorphism can be a unit structure.
The decomposition asserted here is only over $R$, not a split
decomposition of $C$-modules.

There is a natural surjection $\theta_i:F_C^*L_i\to L_i$,
since $(x_i^p)\subset(x_i)$, where $x_1=x$ and $x_2=x-\Delta$.
It linearizes the branch map $r\mapsto r^p$. Its precise kernels are
\[
 \ker\theta_1\simeq L_2/\Delta^{p-1}L_2,
 \qquad \ker\theta_2\simeq L_1/\Delta^{p-1}L_1:
\]
the first is generated by the class of $x$, annihilated by $x-\Delta$,
and the second by the class of $x-\Delta$, annihilated by $x$.
For the $p^e$-power Frobenius all formulas replace $p-1$ by
$p^e-1$; in particular the coefficient-field-linear $q$th-power
version has length $m(q-1)$ per branch.

\par\smallskip\noindent\textbf{Derived pullback on the same object.}\quad
Put $f_1=x$, $f_2=x-\Delta$. Cancellation in the domain $R[X]$
gives $\operatorname{ann}_C(f_1)=(f_2)$ and
$\operatorname{ann}_C(f_2)=(f_1)$, so
\[
 \cdots\xrightarrow{f_1}C\xrightarrow{f_2}C
       \xrightarrow{f_1}C\longrightarrow L_1\longrightarrow0
\]
is an exact free resolution. Frobenius pullback replaces its maps by
$\Delta^{p-1}f_1$ and $\Delta^{p-1}f_2$.
The scalar $\Delta$ is a nonzerodivisor on the $R$-free ring $C$,
so their kernels are still $(f_2)$ and $(f_1)$, while the preceding
images are only $\Delta^{p-1}(f_2)$ and $\Delta^{p-1}(f_1)$.
Writing $C_{F_C}$ for the right module twisted by Frobenius, we obtain
\[
 \operatorname{Tor}^{C}_{2j+1}(C_{F_C},L_1)
     \simeq L_1/\Delta^{p-1}L_1\quad(j\ge0),
\]
\[
 \operatorname{Tor}^{C}_{2j}(C_{F_C},L_1)
     \simeq L_2/\Delta^{p-1}L_2\quad(j\ge1).
\]
Exchange the branches for $L_2$. Every positive degree has
underlying $R$-module $R/(\Delta^{p-1})$.
Thus $\mathbf L F_C^*L_i$ is unbounded below and cannot equal the
degree-zero branch object. Changing its resolution cannot repair
this; replacing the object or using an unbounded category is a
different construction.

\par\smallskip\noindent\textbf{The actual algebra action and multiplicities.}\quad
The completed good-characteristic model from \fid{F19} has
\[
 A=R\times R\times M_2(C).
\]
Define $\phi_A$ by $F_R$ on the scalar blocks and entrywise $F_C$
on the matrix block. It is a ring homomorphism because, in the
commutative coefficient ring $C$,
\[
 \left(\sum_j a_{ij}b_{j\ell}\right)^p
       =\sum_j a_{ij}^p b_{j\ell}^p.
\]
This is not the noncommutative matrix-power map of \fid{F11}.
It fixes the matrix units, so its extension-of-scalars functor
commutes with the standard Morita equivalence.

The actual right module is
\[
 M=R^{20}\oplus R^4\oplus
    (L_1\otimes_C C^{1\times2})^{16}
       \oplus(L_2\otimes_C C^{1\times2})^4,
\]
through the regular action-preserving basis of \fid{F21}. Under
$F_A^*M=M\otimes_{A,\phi_A}A$, each branch multiplicity module
becomes $F_C^*L_i$; the two scalar sectors stay free. Therefore
\begin{equation}\label{eq:frobtorsion}
 F_A^*M\simeq R^{64}\oplus
          (R/(\Delta^{p-1}))^{40}
 \quad\text{as }R\text{-modules}.
\end{equation}
The branch multiplicity sheaf has 20 copies and torsion length
$20m(p-1)$; the actual row module has 40 copies and length
$40m(p-1)$. The natural algebra-linearized surjection discards
exactly this torsion. Derived Morita equivalence transfers the
resolution above as well: every positive Tor degree has 40 copies
of that torsion module, and the scalar sectors contribute none.

For the original center take $t=v-1$, $x=z_1-z$ and
$\Delta=z_1-z_2$. Its formula is
\[
 \Delta=\frac{(D-1)(D-4)}{D^2}
        =\frac34t^2+O(t^3).
\]
In characteristic greater than three its order is exactly two,
giving the claimed length $80(p-1)$.

\par\smallskip\noindent\textbf{Parameter pullback and the regular unit repair.}\quad
Parameter-base pullback of the center gives
\[
 C^{(p)}=R\otimes_{F_R,R}C
       \simeq R[X]/(X(X-\Delta^p)).
\]
Relative Frobenius sends $X$ to $x^p=\Delta^{p-1}x$.
In the two free bases its matrix is
$\operatorname{diag}(1,\Delta^{p-1})$; it is injective with
cokernel of length $m(p-1)$. Its source has branch separation of
order $mp$, whereas the target has order $m$. Extension over this
relative map is the additional operation in center pullback which
causes the torsion.

In the adapted free $R$-basis of $M$, take coefficientwise
Frobenius on every scalar and branch row. Branch evaluation
$C\to R$ commutes with Frobenius, so this semilinear map satisfies
\[
 \tau_{\rm ad}(ma)=\tau_{\rm ad}(m)\phi_A(a).
\]
Its parameter linearization
$R\otimes_{F_R,R}M\to M$ is an isomorphism, sending each
$1\otimes e_j$ to $e_j$. This tensor product is different from
the algebra tensor product in \eqref{eq:frobtorsion}.

The transport to the original tensor basis is explicit. With the
right-row convention of \fid{F21}, write
\[
 \rho_{\rm ad}(a)=B\rho_{\rm old}(a)B^{-1},\qquad
 r_{\rm old}=r_{\rm ad}B,\qquad
 U=(B^{[p]})^{-1}B,
\]
where $B^{[p]}$ means entrywise Frobenius and $B$ is the regular
invertible adapted basis matrix. Then
\begin{equation}\label{eq:frobtransport}
 \tau_{\rm old}(r)=r^{[p]}U,\qquad
 \rho_{\rm old}(\phi_A(a))
       =U^{-1}\rho_{\rm old}(a)^{[p]}U.
\end{equation}
Indeed transform $r$ to $rB^{-1}$, apply entrywise Frobenius,
and multiply back by $B$; applying the same operation to the
action matrix proves the second equality. It also directly checks
$\tau_{\rm old}(r\rho(a))=\tau_{\rm old}(r)\rho(\phi_A(a))$.
The matrix $U$ is invertible, proving the unit parameter structure
on the actual lattice.

Naive entrywise Frobenius in the original tensor basis is not
assumed to preserve the original algebra image. Although its
generator entries have the numerical form $P_i(v)^{[p]}=P_i(v^p)$,
the intrinsic algebra map acts by the conjugated expression in
\eqref{eq:frobtransport}. The regular conjugation cannot be omitted
without an additional argument.

\par\smallskip\noindent\textbf{A unit model on the normalization.}\quad
The normalization is
\[
 C\hookrightarrow\bar C=R\times R,
 \qquad a+bx\longmapsto(a,a+b\Delta).
\]
Its image is the pairs congruent modulo $\Delta$, and its quotient
has length $m$. It is the full normalization: the branch idempotents
are integral, and each factor $R$ is integrally closed.
Extending the algebra to this normalized center gives
\[
 A\hookrightarrow\bar A
       =R\times R\times M_2(R)\times M_2(R),
 \qquad \operatorname{length}_R(\bar A/A)=4m.
\]
Let $f:\bar S\to S$ normalize the three-component scheme in
\fid{F21}; it has four components upstairs. Choose free bundles
\[
 \bar E=(R,R,R^2,R^2)=f^*E,\qquad
 \bar{\mathcal F}=(R^{20},R^4,R^{16},R^4).
\]
They satisfy exactly
\[
 f_*\bar{\mathcal F}=\mathcal F,\qquad
 f_*\mathcal Hom(f^*E,\bar{\mathcal F})
       \simeq\mathcal Hom(E,\mathcal F).
\]
The last is ordinary tensor--Hom adjunction for the finite locally
free $E$. Its right $A$ action is restriction of the $\bar A$
action along the displayed inclusion. It therefore realizes the
same actual module, not just the same dimension or character.

On the normalized free branches, the chosen bases give unit
coherent Frobenius structures on $\bar E$ and
$\bar{\mathcal F}$ and hence the corresponding row module.
Entrywise Frobenius on $\bar A$ commutes with its inclusion of
$A$, since branch evaluations commute with absolute Frobenius.
But the natural comparison
\[
 F_S^*f_*\bar{\mathcal F}
       \longrightarrow f_*F_{\bar S}^*\bar{\mathcal F}
       \simeq f_*\bar{\mathcal F}
\]
is, on each branch, precisely $C/(x_i^p)\twoheadrightarrow C/(x_i)$.
It has torsion length $20m(p-1)$ on the multiplicity sheaf and
$40m(p-1)$ on the actual row module. Thus the unit upstairs
becomes the computed nonunit map downstairs. After inverting $t$
the normalization is an isomorphism and the entire defect vanishes.
\end{proof}

\paragraph{The Grothendieck dictionary and the exact remaining comparison.}
This is a calculation of finite maps, coherent modules, relative
Frobenius and derived tensor products on an explicit singular scheme.
The relative map from collision order $mp$ to order $m$ is a finite
universal homeomorphism, like those in \fid{F22}; its coherent cokernel
is nevertheless nonzero. Ordinary constant \'etale traces lose this
length. The normalization has unit coherent Frobenius and preserves
the original algebra action after restriction, while pushforward
through the collision does not preserve the unit property.

Thus the failed claim is a unit center/algebra pullback on the
specified branch object; the successful claims are the unit
parameter-base structure and the action-compatible normalized model.
Replacing the object, permitting unbounded constructions, or passing
to an \'etale/Weil realization are separate operations. None here
identifies quantum degree with a Weil weight, chooses VIII's basis,
or proves the opposite path-output sign needed by the main criterion.
